\begin{explanationbox}
\textbf{\textcolor{CExplanation}{一句话直觉}}
一条弦的两个端点关于中点 $M$ 对称。椭圆方程里只有平方项；\textbf{把对称两点代进去再相减，平方项自动变成“一加一减”的一次关系}，弦的方向就出来了。
\medskip
\textbf{\textcolor{CExplanation}{最小认知负担：只设一个位移参数}}
当斜率 $k$ 存在时，让两端点沿弦从中点等距离移动：
\[
 A(x_0+t,y_0+kt),\qquad B(x_0-t,y_0-kt),\qquad t\ne0.
\]
横坐标一个加 $t$、一个减 $t$；沿同一条斜率为 $k$ 的直线走，纵坐标自然加减 $kt$。代入椭圆作差，得到
\[
 4t\left(\frac{x_0}{a^2}+\frac{k y_0}{b^2}\right)=0
 \quad\Longrightarrow\quad \frac{x_0}{a^2}+\frac{k y_0}{b^2}=0.
\]
\textbf{理解关键：}不是凭空想到某个公式，而是“中点对称 $\to$ 正负同一个位移 $\to$ 平方差消去 $t^2$”。
\medskip
\textbf{\textcolor{CExplanation}{更深一层：斜率只是无除法关系的特例}}
设端点差为 $\Delta x=x_1-x_2,\ \Delta y=y_1-y_2$。点差法首先得到
\[
 \boxed{\frac{x_0}{a^2}\Delta x+\frac{y_0}{b^2}\Delta y=0.}
\]
这表明弦的方向向量 $(\Delta x,\Delta y)$ 与 $(x_0/a^2,y_0/b^2)$ 垂直。\textbf{先得方向的垂直关系，再在 $\Delta x\ne0$ 时求斜率}，不容易漏掉竖直弦。
\medskip
\begin{center}\begin{tikzpicture}[x=1.45cm,y=1.45cm,>=stealth,font=\small]
\node[font=\small\bfseries] at (0,3.15) {真实弦：中点在内部；弦线只与小椭圆相切};
\draw[->,gray] (-3.5,0)--(3.7,0) node[right] {$x$};
\draw[->,gray] (0,-2.4)--(0,2.5) node[above] {$y$};
\draw[CStatement,thick] (0,0) ellipse (3 and 2);
\coordinate (A) at (3,0); \coordinate (B) at (0,2); \coordinate (M) at (1.5,1);
\draw[teal,very thick] (A)--(B);
\fill (A) circle(1.8pt) node[below right] {$A$};
\fill (B) circle(1.8pt) node[above left] {$B$};
\fill[orange!80!black] (M) circle(2pt) node[above right] {$M$};
\node[below left] at (0,0) {$O$};
\draw[purple,dashed,thick] (0,0) ellipse (2.121320344 and 1.414213562);
\draw[orange!80!black,dashed] (0,0)--(M);
\node[align=center] at (-.8,-.8) {同心小椭圆\\$r^2=1/2$};
\node[below,align=center] at (0,-2.85) {$M\pm s(1,-2/3),\ s=3/2$：两交点关于 $M$ 对称};
\end{tikzpicture}
\end{center}
图中弦的两端点关于 $M$ 对称；它与同心小椭圆在 $M$ 处相切。$M$ 在原椭圆内部，不能当作原椭圆的切点。椭圆伸缩成圆的几何解释见证明四。
\medskip
\textbf{\textcolor{CExplanation}{识别信号和解题顺序}}
见到“\textbf{椭圆 + 弦中点}”，先查中点是否在内部，再写\textbf{无除法弦方程}；只有题目需要斜率且 $y_0\ne0$，才把它化为 $k=-b^2x_0/(a^2y_0)$。需要求中点时，再加“中点在已知直线上”的条件联立。
\end{explanationbox}
